\documentclass[11pt]{article}

% ===================== PACKAGES =====================
\usepackage[a4paper,margin=1in]{geometry}
\usepackage{amsmath,amssymb}
\usepackage{enumitem}
\usepackage{fancyhdr}
\usepackage{booktabs} 
\usepackage{xcolor} 

% ===================== HEADER & FOOTER =====================
\pagestyle{fancy}
\fancyhf{}
\lhead{CSE 400: Fundamentals of Probability in Computing}
\rhead{Tutorial-1 Scribe}
\cfoot{\thepage}

% ===================== DOCUMENT START =====================
\begin{document}

% --- TITLE BLOCK ---
\begin{center}
    {\Large \textbf{Lecture Scribe}} \\
    \vspace{0.3cm}
    \begin{tabular}{rl}
        \textbf{Course:} & CSE 400 – Fundamentals of Probability in Computing \\
        \textbf{Lecture Type:} & Tutorial Session (Problem Solving) \\
        \textbf{Topic:} & Tutorial-1 Solutions \\
        \textbf{Date:} & February 13, 2026 \\
        \textbf{Scribe:} & Het Patel
    \end{tabular}
\end{center}
\hrule
\vspace{0.5cm}

% ===================== CONTENT =====================

\section{Opening Context}
This lecture was a tutorial session in which all assigned problems were solved sequentially. \\
The focus was:
\begin{itemize}
    \item Counting techniques (multinomial coefficients)
    \item Uniform probability modeling
    \item Complement rule and De Morgan’s laws
    \item Poisson distribution
    \item Conditional probability
    \item Gaussian distribution and Q-function
    \item Binomial systems and reliability comparison
\end{itemize}
Each question was solved step-by-step during class.

\section{Q1. Division into Platters (Multinomial Counting)}

\subsection{Problem Statement}
20 distinct dishes are to be divided into 4 groups of 5 dishes each.
\begin{itemize}
    \item The order of dishes within groups does not matter.
    \item The order of groups does not matter.
\end{itemize}

\subsection{Q1(a) Total Number of Ways}
\textbf{Assumptions}
\begin{itemize}
    \item 20 distinct objects
    \item Groups are unlabeled
    \item Each group contains exactly 5 dishes
    \item Order inside group irrelevant
\end{itemize}

\textbf{Result Used} \\
Multinomial coefficient for partitioning into equal groups:
\[ \frac{20!}{(5!)^4 4!} \]

\textbf{Explanation of Terms}
\begin{itemize}
    \item $20!$: Permute all dishes
    \item $(5!)^4$: Remove ordering inside each group
    \item $4!$: Remove ordering between groups
\end{itemize}

\subsection{Q1(b) Probability Each Platter Contains Single Cuisine}
\textbf{Given}
\begin{itemize}
    \item Exactly 5 dishes per cuisine
    \item 4 cuisines
    \item Platters unlabeled
\end{itemize}

\textbf{Key Observation} \\
Only one arrangement satisfies this condition: Each cuisine forms exactly one platter.

\textbf{Probability}
\[ P = \frac{1}{\frac{20!}{(5!)^4 4!}} = \frac{(5!)^4 4!}{20!} \]

\section{Q2. Bus Arrival – Uniform Distribution}

\subsection{Model}
\begin{itemize}
    \item Passenger arrival time uniformly distributed over 30 minutes (7:00–7:30).
    \item Buses every 15 minutes $\rightarrow$ at 7:00, 7:15, 7:30.
\end{itemize}

\subsection{Q2(a) Wait < 5 Minutes}
Waiting < 5 minutes occurs if arrival is:
\begin{itemize}
    \item 7:10–7:15
    \item 7:25–7:30
\end{itemize}
Total favorable time = 10 minutes
\[ P = \frac{10}{30} = \frac{1}{3} \]

\subsection{Q2(b) Wait > 10 Minutes}
Arrival intervals:
\begin{itemize}
    \item 7:00–7:05
    \item 7:15–7:20
\end{itemize}
Total favorable time = 10 minutes
\[ P = \frac{10}{30} = \frac{1}{3} \]

\section{Q3. Conditional Probability Using Complements}

\subsection{Given}
\begin{itemize}
    \item $P(L) = 0.15$
    \item $P(E) = 0.25$
    \item $P(L \cap E) = 0.08$
\end{itemize}
Where:
\begin{itemize}
    \item $L$: arrives late
    \item $E$: leaves early
\end{itemize}
We need:
\[ P(L^c \mid E^c) \]

\textbf{Step 1: Compute Complements}
\[ P(E^c) = 1 - 0.25 = 0.75 \]

\textbf{Step 2: Compute $P(L \cup E)$}
\begin{align*}
    P(L \cup E) &= P(L) + P(E) - P(L \cap E) \\
    &= 0.15 + 0.25 - 0.08 \\
    &= 0.32
\end{align*}

\textbf{Step 3: Use De Morgan’s Law}
\begin{align*}
    P(L^c \cap E^c) &= 1 - P(L \cup E) \\
    &= 1 - 0.32 \\
    &= 0.68
\end{align*}

\textbf{Step 4: Apply Conditional Probability Definition}
\begin{align*}
    P(L^c \mid E^c) &= \frac{P(L^c \cap E^c)}{P(E^c)} \\
    &= \frac{0.68}{0.75} \\
    &\approx 0.907
\end{align*}

\section{Q4. Poisson Distribution ($\lambda = 0.2$)}

\subsection{Model}
$X \sim \text{Poisson}(\lambda = 0.2)$ \\
PMF:
\[ P(X=k) = e^{-\lambda} \frac{\lambda^k}{k!} \]

\subsection{Q4(a) $P(X = 0)$}
\[ P(0) = e^{-0.2} = 0.8187 \]

\subsection{Q4(b) $P(X \ge 2)$}
Using complement:
\begin{align*}
    P(X \ge 2) &= 1 - P(0) - P(1) \\
    P(1) &= 0.2 e^{-0.2} = 0.1637 \\
    P(X \ge 2) &= 1 - 0.8187 - 0.1637 \\
    &= 0.0176
\end{align*}

\section{Q5. Poisson ($\lambda = 3.5$)}
$X \sim \text{Poisson}(3.5)$

\subsection{Q5(a) $P(X \ge 2)$}
\begin{align*}
    P(X \ge 2) &= 1 - P(0) - P(1) \\
    P(0) &= e^{-3.5} = 0.0302 \\
    P(1) &= 3.5e^{-3.5} = 0.1057 \\
    P(X \ge 2) &= 1 - 0.0302 - 0.1057 \\
    &= 0.8641
\end{align*}

\subsection{Q5(b) $P(X \le 1)$}
\begin{align*}
    P(X \le 1) &= P(0) + P(1) \\
    &= 0.0302 + 0.1057 \\
    &= 0.1359
\end{align*}

\section{Q6. Conditional Probability in Poisson ($\lambda = 3$)}
$X \sim \text{Poisson}(3)$

\subsection{Q6(a) $P(X \ge 3)$}
\begin{align*}
    P(X \ge 3) &= 1 - P(0) - P(1) - P(2) \\
    &= 1 - (0.0498 + 0.1494 + 0.2240) \\
    &= 0.5768
\end{align*}

\subsection{Q6(b) $P(X \ge 3 \mid X \ge 1)$}
Definition:
\[ P(A \mid B) = \frac{P(A \cap B)}{P(B)} \]
Since $X \ge 3 \subseteq X \ge 1$:
\begin{align*}
    P(X \ge 3 \mid X \ge 1) &= \frac{P(X \ge 3)}{P(X \ge 1)} \\
    P(X \ge 1) &= 1 - P(0) = 0.9502 \\
    P(X \ge 3 \mid X \ge 1) &= \frac{0.5768}{0.9502} \approx 0.607
\end{align*}

\section{Q7. Gaussian Random Variable}

\subsection{Given}
\[ f_X(x) = \frac{1}{\sqrt{8\pi}} \exp\left(-\frac{(x+3)^2}{8}\right) \]

\subsection{Parameter Identification}
Comparing with standard Gaussian:
\begin{itemize}
    \item $m = -3$
    \item $\sigma^2 = 4 \implies \sigma = 2$
\end{itemize}

\subsection{Results}
\begin{align*}
    P(X \le 0) &= \Phi(3/2) = 1 - Q(3/2) \\
    P(X > 4) &= Q(7/2) \\
    P(|X+3| < 2) &= 1 - 2Q(1) \\
    P(|X-2| > 1) &= 1 - Q(2) + Q(3)
\end{align*}

\section{Q8. Jury Decision Problem}

\subsection{Let:}
\begin{itemize}
    \item 12 jurors
    \item Conviction requires $\ge 8$ guilty votes
    \item Each juror correct with probability $\theta$
    \item $\alpha$ = probability defendant guilty
\end{itemize}

\textbf{If Guilty:}
\[ \sum_{i=8}^{12} \binom{12}{i} \theta^i (1-\theta)^{12-i} \]

\textbf{If Innocent:}
\[ \sum_{i=5}^{12} \binom{12}{i} \theta^i (1-\theta)^{12-i} \]

\textbf{Final Expression (Conditioning):}
\[ \alpha \sum_{i=8}^{12} \binom{12}{i} \theta^i (1-\theta)^{12-i} + (1-\alpha) \sum_{i=5}^{12} \binom{12}{i} \theta^i (1-\theta)^{12-i} \]

\section{Q9. Poisson Form Identification}

\subsection{Given}
\[ p(i) = c \frac{\lambda^i}{i!} \]
Using:
\[ \sum_{i=0}^{\infty} \frac{\lambda^i}{i!} = e^\lambda \]
Thus:
\[ c e^\lambda = 1 \implies c = e^{-\lambda} \]

\subsection{Results}
\begin{align*}
    P(X=0) &= e^{-\lambda} \\
    P(X>2) &= 1 - e^{-\lambda} - \lambda e^{-\lambda} - \frac{\lambda^2}{2} e^{-\lambda}
\end{align*}

\section{Q10. System Reliability}

\subsection{Model}
Number functioning:
\[ X \sim \text{Binomial}(n,p) \]

\subsection{Q10(a)}
5-component system better than 3-component system if:
\[ p > \frac{1}{2} \]

\subsection{Q10(b)}
General result: \\
A $2k+1$-component system is better than $2k-1$-component system
\[ \iff p > \frac{1}{2} \]

Derived expression:
\[ \binom{2k-1}{k} p^k (1-p)^k (2p-1) \]
Sign depends only on $2p-1$.

\section{Final Summary of Tutorial}
Concepts used in this tutorial:
\begin{itemize}
    \item Multinomial coefficients
    \item Uniform distribution
    \item Complement rule
    \item De Morgan’s laws
    \item Poisson distribution
    \item Conditional probability
    \item Gaussian distribution
    \item Q-function
    \item Binomial reliability comparison
\end{itemize}
All results were derived through direct substitution and algebraic simplification as demonstrated in class.

\vspace{0.5cm}
\begin{center}
    \small \textit{End of Scribe}
\end{center}

\end{document}